\section{Finite norm and expectation estimates}\label{sec:moments}
\subsection{Projection and finite graded norm bounds}
Let $\widetilde K=K\otimes I_3$ and use the finite occupation creators to define
$\mathcal S_B=\tfrac12\sum_{ij}B_{ij}C_iC_j$ and
$\Psi_B=\sum_{j=0}^{A/2}\mathcal S_B^j\vac/j!$.
The one-sided identity
\begin{equation}\label{eq:projection}
 J_{\rm hc}^*C_i=x_iJ_{\rm hc}^*,\qquad
 \psi_K=J_{\rm hc}^*\Psi_{\widetilde K}
\end{equation}
holds exactly. Creation cannot repair a configuration with two particles on a
rung; hence the one-sided identity holds even on the orthogonal complement of
the hard-core image. It does not assert multiplicativity of arbitrary
compressions. For creation polynomials it gives the needed action intertwiner
and a genuine Euclidean norm contraction.

\begin{lemma}[Finite graded majorant]\label{lem:norm}
Suppose $B$ is a real symmetric mode matrix, $B^2\le qI$, and
$\tr B^2\le T$, with $q,T\ge0$. Define
\begin{equation}\label{eq:bm}
 b_0=1,\qquad b_{j+1}=b_j\frac{qj+T/2}{j+1},\qquad
 M_{A,d}=\sum_{j=0}^{A/2}\binom{2j+d}{d}b_j.
\end{equation}
Then $\norm{\Psi_B}^2\le M_{A,0}$. For a homogeneous creation polynomial
$P_d$, with faithful vacuum norm $\norm{P_d\vac}^2\le F_d$, one has
\begin{equation}\label{eq:amplify}
 \norm{P_d\Psi_B}^2\le F_d M_{A,d}.
\end{equation}
\end{lemma}
\begin{proof}
On a safe grade the exact pair commutator is
$[\mathcal S_B^*,\mathcal S_B]=d\Gamma(B^2)+\tr(B^2)/2$.
At the top two grades the missing creations contribute with the favorable
nonpositive boundary correction. If $A_{2j-2}$ bounds the squared creation
norm on grade $2j-2$, finite adjoint duality gives
$\norm{\mathcal S_B^*\xi_{2j}}^2\le A_{2j-2}\norm{\xi_{2j}}^2$.
The commutator therefore yields $A_{2j}\le A_{2j-2}+2qj+T/2$, starting
from $A_0=T/2$. Summing this finite recurrence gives
\[
 \norm{\mathcal S_B\xi_{2j}}^2
       \le(j+1)(qj+T/2)\norm{\xi_{2j}}^2.
\]
Division by factorials yields \eqref{eq:bm}; distinct grades are orthogonal.
The finite symmetric-tensor multiplication estimate on grade $n$ is
$\norm{P_d\xi_n}^2\le\binom{n+d}{d}\norm{P_d\vac}^2\norm{\xi_n}^2$.
Here is an occupation-basis proof of that multiplication estimate. Write
$r_\alpha=(P_d\vac)_\alpha$ for $|\alpha|=d$, and let $v_\nu$ be the
coefficients of $\xi_n$. For $|\gamma|=n+d$ the output coefficient is
\[
 (P_d\xi_n)_\gamma=
 \sum_{\substack{\alpha\le\gamma\\|\alpha|=d}}
 r_\alpha v_{\gamma-\alpha}
 \sqrt{\frac{\gamma!}{\alpha!(\gamma-\alpha)!}}.
\]
The multivariate Vandermonde identity gives
$\sum_{\alpha\le\gamma,|\alpha|=d}\binom\gamma\alpha=\binom{n+d}{d}$.
Cauchy--Schwarz therefore bounds the squared coefficient by
$\binom{n+d}{d}\sum_\alpha|r_\alpha|^2|v_{\gamma-\alpha}|^2$.
Summing over $\gamma$ gives the product of the two squared norms; finite
cutoffs only delete nonnegative output terms. Applying the estimate grade by
grade and summing proves \eqref{eq:amplify}.
The implemented occupation proof includes all cutoff-boundary cases; it never
uses an untruncated CCR at top degree.
\end{proof}

For the cubic kernel the exact absolute row bound is
$\kappa(g)=g+2g^2+13g^3/4\le493/4000$.
The squared row norms at the endpoint, which dominate the full interval, are
\begin{equation}\label{eq:rowSquares}
 \rho_{16}=\frac{362417}{128000000},\quad
 \rho_{32}=\frac{1443579}{512000000},\quad
 \rho_{64}=\frac{143749}{51200000}.
\end{equation}
The endpoint domination is algebraic: on these bipartite graphs,
\[
 |K|=\frac g4D+\frac{g^2}{8}(D_2+D)
       +\frac{g^3}{64}(5D^3+4D_2-40D).
\]
Every displayed coefficient matrix is nonnegative. On an edge $(D^3)_{rs}\ge8$;
on other entries the last assertion is immediate, and bipartite parity
separates $D_2$ from $D,D^3$. Thus each absolute entry and squared row sum is
monotone throughout the interval. Use $q=(493/4000)^2$ and $T_A=3A\rho_A$. All sums in \eqref{eq:bm} are
finite; the largest has 33 terms. They apply both to signed $\widetilde K$ and
to its entrywise absolute value, without equating the two state norms.

\subsection{Faithful residuals and the matching Gram matrix}
For a distinct four-site set $S=\{s_0,s_1,s_2,s_3\}$, let $Q_0,Q_1,Q_2$ be the
products of two color-contracted pair creators corresponding to
$01|23$, $02|13$, $03|12$. Collect the \emph{signed} contributions of every
bond and every repeated description of the same matching into $c_{S,m}(g)$.
Then $R_4=\sum_{S,m}c_{S,m}Q_m$. Terms with repeated sites vanish physically
and are removed before choosing its faithful occupation lift.

\begin{lemma}[Four-site Gram and translation count]\label{lem:gram}
For distinct sites,
\[
 \ip{Q_m\vac}{Q_{m'}\vac}=\begin{cases}9&m=m',\\3&m\ne m',\end{cases}
 \quad
 \norm{\sum_m c_mQ_m\vac}^2
     =6\sum_m c_m^2+3\left(\sum_m c_m\right)^2.
\]
Different four-site supports are orthogonal. For any translation-invariant
support weight $w(S)$ on an $A$-site torus,
\begin{equation}\label{eq:rootcount}
 \sum_{|S|=4}w(S)=\frac A4\sum_{\substack{|S|=4\\0\in S}}w(S).
\end{equation}
\end{lemma}
\begin{proof}
Each matching has nine color assignments. Two different matchings share exactly
the three assignments in which all four colors agree. Distinct supports give
distinct occupation supports. For \eqref{eq:rootcount}, count pairs $(r,S)$ with
$r\in S$ in two ways and translate $r$ to zero. This counts incidences, not
translation orbits; it remains correct for supports with nontrivial stabilizers.
\end{proof}

An explicit local specification makes the signed coefficient collection
independent of the table implementation. For site functions $f,h$ define
$P(f,h)=\sum_\alpha(\sum_i f_ix_{i\alpha})(\sum_j h_jx_{j\alpha})$.
Extend the outside rows $u,v$ by zero at both edge endpoints. For a matching
$m$ of a distinct support $S$, put
\[
 \mathcal T_m(f,h;j,l)=
 \sum_{\substack{(a,b,c,d)\text{ bijective onto }S\\
                 \{\{a,b\},\{c,d\}\}=m}}f_a h_b j_c l_d.
\]
There are exactly eight ordered assignments. Expanding \eqref{eq:q4} gives
\begin{align*}
 q_4={}&-P(e_r,u)P(u,v)-P(e_s,v)P(u,v)+P(e_r,v)P(e_s,u)\\
 &-(1+k)P(e_r,e_s)P(u,v)-P(e_r,u)P(e_s,u)\\
 &-P(e_r,v)P(e_s,v)+3kP(e_r,u)P(e_s,v).
\end{align*}
The coefficient $c_{S,m}$ is exactly $g/2$ times the edge sum of these seven
terms with each product $P(f,h)P(j,l)$ replaced by $\mathcal T_m(f,h;j,l)$.
This explicitly defines the fiber sum before any absolute value.

The pair residual has a symmetric zero-diagonal coefficient $F$ with
$R_2=\frac12\sum_{r,s,\alpha}F_{rs}x_{r\alpha}x_{s\alpha}$ and
$\norm{R_2\vac}^2=(3A/2)\sum_sF_{0s}^2$.
The sextic residual is lifted after deleting repeated-site terms. Define its
unfiltered absolute scalar mass by
\[
 \ell_6(g)=\frac g2\sum_{\{r,s\}\in\mathcal E}
 \left(\sum_{t\ne r,s}|K_{rt}(g)|\right)^2
 \left(\sum_{t\ne r,s}|K_{st}(g)|\right)^2.
\]
Set $\kappa:=493/4000$. Each product of three color-contracted pairs on distinct sites has vacuum norm
$\sqrt{27}$, giving $F_6\le27\ell_6(g)^2$. The entire outside sums are
monotone in $g$, so they may first be moved to $g=1/10$ and only then evaluated
by subtracting the endpoint bond from $\kappa$. The endpoint horizontal and
vertical bond numerators over $64000$ are $(1690,1690)$, $(1685,1690)$, and
$(1685,1685)$ on the three boxes. Hence the exact bound is
\[
 \ell_6(g)\le\frac A{20}
 \left[(\kappa-B_x/64000)^4+(\kappa-B_y/64000)^4\right].
\]
These are precisely the rational masses in Table~\ref{tab:vacuum} and its
64-rung companion in the source. This deliberately conservative bound requires no unjustified
orthogonality between overlapping sextic descriptions.

\subsection{A positive finite Wick upper bound}
For the expectation estimate one must use the positive matrix
$B=|K|\otimes I_3$, not a signed-kernel Gaussian identity.
Let $\Psi_+=\Psi_B$, $G=\norm{\Psi_+}^2$, and
\[
 \tau_A=\frac{\kappa\rho_A}{1-\kappa^2},\qquad
 L_{rs}=|K_{rs}|+\tau_A,\qquad \mathcal L=L\otimes I_3,
 \quad\kappa=493/4000.
\]
\begin{lemma}[Finite Wick domination]\label{lem:wick}
For every even list of mode indices $I=(i_1,\ldots,i_{2d})$,
\begin{equation}\label{eq:wick}
 \ip{\Psi_+}{C_{i_1}\cdots C_{i_{2d}}\Psi_+}
 \le G\sum_{\pi}\prod_{\{a,b\}\in\pi}\mathcal L_{i_a i_b},
\end{equation}
where $\pi$ runs through the perfect matchings of the list positions.
\end{lemma}
\begin{proof}
All occupation coefficients are nonnegative. The exact truncated commutator
satisfies $C_i^*C_j\le C_jC_i^*+\delta_{ij}I$ entrywise. Expanding the finite
Taylor polynomial gives $C_i^*\Psi_+\le\sum_j B_{ij}C_j\Psi_+$.
Moving one annihilator through a creation word and then using this inequality
again yields, for an odd list $I$, the actual recurrence
\[
 f_i\le\sum_{k\in I}B_{i,i_k}F_{I\setminus k}+\sum_j(B^2)_{ij}f_j,
 \quad f_i=\ip{\Psi_+}{C_i C_I\Psi_+}.
\]
The sum is over occurrences, so repeated indices are counted correctly.
For the rung block, symmetry and the squared-row bound imply
$(|K|^2)_{ij}\le\rho_A$ and $(|K|^3)_{ij}\le\kappa\rho_A$.
Consequently $L\ge|K|+|K|^2L$ entrywise and the row sums of $|K|^2$ are at most
$\kappa^2<1$. If a finite vector satisfies $h\le B^2h$, its positive maximum
would satisfy $h_i\le\kappa^2h_i<h_i$; hence $h\le0$.
Induction on the word length, comparing $f$ with the perfect-matching expression,
now proves \eqref{eq:wick}. The empty-word moment is exactly $G$.
This is an upper bound for a finite Taylor state, not an exact Gaussian Wick
identity. Color blocks remain separate throughout.
\end{proof}

Coefficientwise absolute values and \eqref{eq:projection} give
$|\psi_K|\le J_{\rm hc}^*\Psi_+$. Replacing a signed creation polynomial by
its coefficientwise absolute majorant and deleting the hard-core projection
can only increase the resulting nonnegative scalar product. Thus
Lemma~\ref{lem:wick} bounds the \emph{actual physical signed residual mean}.
For four sites put $p_m=\prod_{\{a,b\}\in m}L_{s_as_b}$.
Summing the two free colors gives the weight
\begin{equation}\label{eq:quarticweight}
 9p_m+3\sum_{m'\ne m}p_{m'}=6p_m+3\sum_{m'}p_{m'}.
\end{equation}
For six sites, $27$ color assignments and $15$ positional matchings give the
upper bound $405\max_{rs}L_{rs}^3$ per unit scalar mass.
All these mean estimates retain the multiplier $G\le M_{A,0}$. No equality of
the signed and positive state norms is asserted.
